\documentclass[11pt,a4paper]{article}

\usepackage[utf8]{inputenc}

\usepackage[T1]{fontenc}

\usepackage[polish]{babel}

\usepackage{lmodern}

\usepackage{geometry}

\geometry{margin=27mm}

\usepackage{amsmath}
\usepackage[hidelinks]{hyperref}

\title{Pewny ton pierwszej odpowiedzi LLM a bezkrytyczne przyjmowanie rekomendacji}

\author{Mateusz Pomianek}

\date{Szkic naukowy, wersja 0.2.1}

\begin{document}
\maketitle
\begin{abstract}
Celem prospektywnego eksperymentu jest ustalenie, czy językowa pewność rekomendacji AI zwiększa podatność na błędną poradę oraz zmniejsza selektywność korzystania z poprawnych i błędnych rekomendacji. Projekt zakłada randomizację tonu i poprawności porady, pomiar decyzji przed poradą i po niej oraz analizę logistyczną z efektami losowymi uczestnika i zadania.
\end{abstract}
\section{Wprowadzenie}
Zaufanie deklarowane i behawioralne poleganie na poradzie stanowią odrębne konstrukty. Przyjęcie trafnej rekomendacji może poprawić decyzję, natomiast przyjęcie błędnej porady może ją pogorszyć. W badaniu rozdzielamy jakość rekomendacji od jej językowej kategoryczności. Ton nie jest miarą skalibrowanego prawdopodobieństwa poprawności modelu. Istniejące badania nad nadmiernym poleganiem na AI \cite{overreliance} stanowią punkt odniesienia; proponowany wkład dotyczy zmian decyzji względem indywidualnego osądu sprzed ekspozycji.
\section{Hipotezy}
H1: wśród prób z poprawną decyzją wstępną błędna rekomendacja przedstawiona zdecydowanie częściej prowadzi do błędnej decyzji końcowej niż identyczna rekomendacja przedstawiona ostrożnie.

H2: w próbach początkowej niezgodności z rekomendacją zdecydowany ton zmniejsza różnicę prawdopodobieństw przyjęcia poprawnej i błędnej porady.

H3 (eksploracyjna): pewność własnego osądu przed ekspozycją moderuje wpływ tonu. Nie utożsamiamy deklarowanej pewności z kompetencją metapoznawczą.
\section{Materiał i metody}
\subsection{Schemat i uczestnicy}
Badanie internetowe osób pełnoletnich, z powtarzanymi próbami w układzie czynnikowym $2\times2$. Ton (ostrożny lub zdecydowany) i poprawność rekomendacji (poprawna lub błędna) będą randomizowane po zapisaniu decyzji wstępnej. Każdy uczestnik otrzyma zadania z jednoznacznym kluczem binarnym, bez powtórzeń treści. Warunki będą równoważone między zadaniami i uczestnikami. Kolejność zadań będzie losowa.
\subsection{Bodźce i procedura}
Każde zadanie obejmuje decyzję wstępną, ocenę pewności w skali 0--100, ekspozycję na jedną zamrożoną rekomendację i decyzję końcową z oceną pewności. Warianty tonu w obrębie tej samej rekomendacji zachowują argumenty, fakty, zalecenie i możliwie podobną długość. Nie stosuje się generowania na żywo. Niezależni oceniający sprawdzają klucz poprawności oraz równoważność treści. Pilotaż oceni postrzeganą pewność, uprzejmość i autorytatywność; kryteria akceptacji bodźców zostaną ustalone przed pilotażem. Uczestnicy pilotażu nie wejdą do próby głównej.
\subsection{Zmienne i estymandy}
Pierwszorzędowy wskaźnik: przejście od poprawnej decyzji wstępnej do błędnej decyzji końcowej w próbach z błędną poradą. Estymandem jest marginalna różnica ryzyka między tonem zdecydowanym i ostrożnym w tej populacji prób. Drugorzędowe wyniki obejmują korektę błędnej decyzji po poprawnej poradzie, selektywność akceptacji przy początkowej niezgodności, trafność końcową oraz kalibrację pewności. Prób zgodności wstępnej nie klasyfikuje się jako dowodu przyjęcia porady.
\section{Plan analizy}
Dla pierwszorzędowego wyniku stosujemy model
\[
Y_{ij}\sim\operatorname{Bernoulli}(p_{ij}),\quad
\operatorname{logit}(p_{ij})=\alpha+\beta T_{ij}+u_i+v_j,
\]
gdzie $T=1$ oznacza ton zdecydowany, a $u_i,v_j$ są efektami losowymi uczestnika i zadania. Raportujemy standaryzowaną różnicę przewidywanych ryzyk z 95-procentowym przedziałem ufności oraz iloraz szans. Dla H2, w próbach początkowej niezgodności, model przyjęcia porady obejmie ton, poprawność i ich interakcję. Kontrast marginalny różnic ryzyk, a nie sam współczynnik interakcji na skali logitowej, będzie miarą selektywności. H1 jest jedynym testem potwierdzającym na poziomie 0,05; H2 i H3 mają status drugorzędowy i eksploracyjny.
\subsection{Liczebność i jakość danych}
Liczebność zostanie ustalona symulacyjnie przed rekrutacją, przy uwzględnieniu minimalnego istotnego efektu, wariancji efektów losowych, liczby zadań i odsetka poprawnych decyzji wstępnych. Nie ustala się liczebności na podstawie istotności uzyskanej w trakcie badania. Prerejestracja określi reguły wykluczeń technicznych, obsługę braków odpowiedzi i procedurę analizy wrażliwości. Nie wykluczamy uczestników ze względu na zgodność z hipotezą. Raportujemy liczebność na każdym etapie i braki według warunków.
\section{Wyniki}
Badanie pozostaje na etapie projektowania; akwizycja danych nie została rozpoczęta.
\section{Dyskusja i ograniczenia}
Zamrożenie estymandu, bodźców i planu analizy nastąpi przed rekrutacją. Przed rekrutacją wymagane są właściwa ocena etyczna, świadoma zgoda, minimalizacja danych i wyjaśnienie po badaniu, że część porad była celowo błędna. Zadania nie dotyczą decyzji medycznych, prawnych ani finansowych o realnych konsekwencjach. Manipulacja dotyczy językowej formy porady, nie rzeczywistej pewności modelu. Wnioskowanie ogranicza się do krótkotrwałych decyzji w kontrolowanych zadaniach; nie dowodzi trwałej utraty autonomii. Efekt pierwszej porady w zadaniu nie jest tożsamy z efektem pierwszego kontaktu z AI.
\section{Wnioski}
Projekt pozwala odróżnić szkodliwe poleganie na błędnej poradzie od korzystnej korekty własnego błędu. Przed rozpoczęciem badania pozostają pilotaż, analiza mocy i prerejestracja.
\begin{thebibliography}{9}
\bibitem{overreliance} \emph{Trust and reliance on AI -- An experimental study on the extent and costs of overreliance on AI}. Computers in Human Behavior (2024). \url{https://doi.org/10.1016/j.chb.2024.108352}.
\end{thebibliography}
\end{document}
